%% ==========================================================================
%%  4  Structure of the Chore Problem
%%
%%  The three reusable tools: the arc-update table (the whole goods/chore
%%  asymmetry), the two-tier characterisation of {0,1} solutions, and the
%%  size-shift transform relating the chore problem to R3's goods problem.
%% ==========================================================================
\section{Structure of the Chore Problem}
\label{sec:structure}

Throughout this section we work in cost form, so
$\edgew{\alloc}(i,j) = \cost_i(\bundle{i}) - \cost_i(\bundle{j})$.

% --------------------------------------------------------------------------
\subsection{The arc-update table, and the whole of the asymmetry}
\label{sec:arcupdate}

Every incremental algorithm in this literature adds one item to one bundle and
then asks what happened to the envy graph. The answer is two lines long, and
comparing those two lines against their goods counterparts accounts for the
entire difference between the present problem and that of~\cite{BKNS22}.

For an allocation $\alloc$, an agent $x$ and an unallocated chore
$g \notin \bigcup_i \bundle{i}$, write
\[
  Z^x \;:=\; (\bundle{1}, \dots, \bundle{x} \cup \set{g}, \dots, \bundle{n})
\]
for the allocation obtained by giving $g$ to $x$, and put
\[
  \delta_x \;:=\; \cost_x(\bundle{x} \cup \set{g}) - \cost_x(\bundle{x}),
  \qquad
  \delta_{i,x} \;:=\; \cost_i(\bundle{x} \cup \set{g}) - \cost_i(\bundle{x}),
\]
both of which lie in $\set{0,1}$. So $\delta_x$ is what the chore costs $x$
herself, and $\delta_{i,x}$ is what it costs in $i$'s opinion.

\begin{lemma}[Arc update]
\label{lem:arcupdate}
For all $i,j \in \agents$,
\[
  \edgew{Z^x}(i,j) = \edgew{\alloc}(i,j) \quad (i,j \ne x),
  \qquad
  \edgew{Z^x}(x,j) = \edgew{\alloc}(x,j) + \delta_x,
  \qquad
  \edgew{Z^x}(i,x) = \edgew{\alloc}(i,x) - \delta_{i,x}.
\]
\end{lemma}

\begin{proof}
Only $\bundle{x}$ changes, so an arc between two agents other than $x$ has both
endpoints untouched. For an arc leaving $x$,
$\edgew{Z^x}(x,j) = \cost_x(\bundle{x} \cup \set{g}) - \cost_x(\bundle{j})
= \edgew{\alloc}(x,j) + \delta_x$. For an arc entering $x$,
$\edgew{Z^x}(i,x) = \cost_i(\bundle{i}) - \cost_i(\bundle{x} \cup \set{g})
= \edgew{\alloc}(i,x) - \delta_{i,x}$.
\end{proof}

The mirror statement for goods is obtained by replacing costs with values: there
the arcs \emph{out of} the receiving agent fall and the arcs \emph{into} her
rise. The consequences are opposite in a way that matters.

\begin{center}
\begin{tabular}{@{}lll@{}}
\toprule
 & goods~\cite{BKNS22} & chores (here) \\
\midrule
arcs out of the receiver & fall & \textbf{rise by $\delta_x$} \\
arcs into the receiver   & \textbf{rise} & fall by $\delta_{i,x}$ \\
so the receiver must be  & \emph{maximally} subsidised & \emph{minimally} subsidised \\
which is guaranteed by   & $\Mset(\subsidy) \ne \emptyset$ & $\subsidy_x = 0$ for some $x$ \\
\bottomrule
\end{tabular}
\end{center}

Both guarantees hold trivially --- some agent is always maximally subsidised,
and some agent always has $\subsidy_x = 0$ under the minimum subsidy vector. The
asymmetry is subtler than the table suggests, and it is what
Theorem~\ref{thm:obstruction} exploits. Under goods, the harm done by an
insertion is confined to paths \emph{ending} at the receiver, and requiring
$x \in \Mset(\subsidy)$ caps every such path at $\subsidy_u - \subsidy_x \le 0$
in one stroke. Under chores, the harm is done to paths \emph{leaving} the
receiver, and $\subsidy_x = 0$ caps only the paths that \emph{start} at $x$. A
path that merely passes through $x$ --- entering along an arc from some $i$ with
$\delta_{i,x} = 0$ and leaving along an arc that has just risen by $\delta_x = 1$
--- is not controlled by any condition on $\subsidy_x$. That gap is not a defect
of the analysis; \Cref{sec:approach1} shows it is realised.

% --------------------------------------------------------------------------
\subsection{The shape of a $\set{0,1}$ solution}
\label{sec:twotier}

\begin{proposition}[Two-tier characterisation]
\label{prop:twotier}
Let $\subsidy \in \set{0,1}^n$ and put $S := \set{i \in \agents \mid \subsidy_i = 1}$.
Then $(\alloc,\subsidy)$ is an envy-free solution if and only if
\begin{enumerate}
  \item[(a)] $\cost_i(\bundle{i}) \le \cost_i(\bundle{j})$ whenever $i,j$ lie
    both in $S$ or both in $\agents \setminus S$;
  \item[(b)] $\cost_i(\bundle{i}) \le \cost_i(\bundle{j}) + 1$ for $i \in S$ and
    $j \in \agents \setminus S$;
  \item[(c)] $\cost_i(\bundle{i}) + 1 \le \cost_i(\bundle{j})$ for
    $i \in \agents \setminus S$ and $j \in S$.
\end{enumerate}
\end{proposition}

\begin{proof}
In cost form Definition~\ref{def:efsolution} reads
$\cost_i(\bundle{i}) \le \cost_i(\bundle{j}) + \subsidy_i - \subsidy_j$ for all
$i,j$. Substituting the three possible values $0, 1, -1$ of
$\subsidy_i - \subsidy_j$ gives (a), (b) and (c) respectively.
\end{proof}

The reading is worth keeping in mind when designing an algorithm. The unpaid
agents are exactly envy-free among themselves and \emph{strictly} prefer their
own bundle to that of every paid agent, by a full unit; the paid agents are
exactly envy-free among themselves and envy-free up to one unit towards the
unpaid. This is a certificate format, and it also suggests choosing the
partition $S \mid \agents \setminus S$ \emph{first} and allocating afterwards ---
an approach not attempted here, and the one place where the incremental template
killed in \Cref{sec:approach1} is not obviously an obstacle. In the proof of
Theorem~\ref{thm:binadd} the set $S$ turns out to be exactly the agents holding
a $\lceil q/n \rceil$-share of the universally disliked chores.

% --------------------------------------------------------------------------
\subsection{The size-shift transform}
\label{sec:sizeshift}

Remark~\ref{rem:complement} showed that complementation relates the chore and
goods models only for $n = 2$. There is a second transform, available for every
$n$, which does not solve the problem but says precisely how far it is from
being a corollary of~\cite{BKNS22}.

\begin{theorem}[Size shift]
\label{thm:sizeshift}
For a negative dichotomous instance with cost functions $\cost_i$, define
$\tilde{v}_i(S) := \abs{S} - \cost_i(S)$. Then each $\tilde{v}_i$ is a
dichotomous \emph{goods} valuation in the sense of~\cite{BKNS22}, and for every
allocation $\alloc$ and all $i,j$,
\[
  \edgew{\alloc}(i,j) \;=\; \tilde{\edgew{}}_{\alloc}(i,j) \;+\;
  \abs{\bundle{i}} - \abs{\bundle{j}},
\]
where $\tilde{\edgew{}}_{\alloc}$ denotes the envy graph of the goods instance
$\set{\tilde{v}_i}$. Consequently every \emph{cycle} has the same weight in both
graphs --- so $\alloc$ is envy-freeable for the chores $\set{\cost_i}$ if and
only if it is envy-freeable for the goods $\set{\tilde{v}_i}$ --- while for a
path $P$ from $u$ to $v$,
\[
  \edgew{\alloc}(P) \;=\; \tilde{\edgew{}}_{\alloc}(P) + \abs{\bundle{u}} -
  \abs{\bundle{v}},
  \qquad\text{hence}\qquad
  \pathw{\alloc}(u) \;=\; \max_{v \in \agents}\Big[\,\tilde{d}_{\alloc}(u,v) +
  \abs{\bundle{u}} - \abs{\bundle{v}}\,\Big],
\]
with $\tilde{d}_{\alloc}(u,v)$ the maximum weight of a path from $u$ to $v$ in
the goods graph and $\tilde{d}_{\alloc}(u,u) = 0$.
\end{theorem}

\begin{proof}
$\tilde{v}_i(\emptyset) = 0$, and for $g \notin S$,
$\tilde{v}_i(S \cup \set{g}) - \tilde{v}_i(S) = 1 - (\cost_i(S \cup \set{g}) -
\cost_i(S)) \in \set{0,1}$, so $\tilde{v}_i$ is dichotomous. For the identity,
\[
  \tilde{\edgew{}}_{\alloc}(i,j) = \tilde{v}_i(\bundle{j}) - \tilde{v}_i(\bundle{i})
  = \big(\abs{\bundle{j}} - \cost_i(\bundle{j})\big) -
    \big(\abs{\bundle{i}} - \cost_i(\bundle{i})\big)
  = \edgew{\alloc}(i,j) - \abs{\bundle{i}} + \abs{\bundle{j}} .
\]
Summing along a path $(i_1,\dots,i_r)$, the terms
$\abs{A_{i_t}} - \abs{A_{i_{t+1}}}$ telescope to
$\abs{A_{i_1}} - \abs{A_{i_r}}$; along a cycle they telescope to $0$. The
characterisation of envy-freeability by the absence of positive cycles
(Theorem~\ref{thm:hs-characterisation}) then transfers, and the formula for
$\pathw{\alloc}(u)$ follows from Theorem~\ref{thm:hs-minsubsidy}.
\end{proof}

\begin{remark}[Why Conjecture~\ref{conj:main} is not a corollary of~\cite{BKNS22}]
\label{rem:notcorollary}
Theorem~\ref{thm:sizeshift} says the chore problem is the goods problem of
\cite{BKNS22} \emph{plus a cardinality-balancing requirement}. A sufficient
condition falls straight out: since $\tilde{d}_{\alloc}(u,v) \le \tilde{p}_u -
\tilde{p}_v$ for any envy-eliminating $\tilde{p}$, an allocation solving the
goods instance in which the quantity $\abs{\bundle{i}} + \tilde{p}_i$ has spread
at most $1$ across agents satisfies $\pathw{\alloc} \le 1$. But~\cite{BKNS22}
offers no control whatever on bundle cardinalities --- its output need not even
be EF1, by its own Appendix~C --- so the requirement is not met by invoking it
as a black box. Nor can its algorithm simply be re-run with a balancing
tie-break, because Theorem~\ref{thm:obstruction} rules out the entire
item-by-item template on which it is built. The transform also shows the chore
problem is at least as expressive as the goods problem, so no easier.
\end{remark}
